\documentclass[11pt]{article}
\usepackage[margin=1in]{geometry}
\usepackage[T1]{fontenc}
\usepackage[utf8]{inputenc}
\usepackage{lmodern}
\usepackage{hyperref}
\usepackage{xcolor}
\usepackage{fancyvrb}
\usepackage{listings}
\usepackage{CJKutf8}
\lstset{basicstyle=\ttfamily\small,columns=fullflexible,breaklines=true,upquote=true,frame=single,framerule=0.2pt,tabsize=2}

\title{\texttt{\detokenize{run_d2co.py}} (Q\&A Documentation)}
\author{}
\date{}
\begin{document}
\begin{CJK*}{UTF8}{gbsn}
\maketitle

\paragraph{Q: What does this run script control?}
\textbf{A:} It controls the full experiment loop: args, data loading, model instantiation, training with a specific loss/optimizer, and evaluation with shared metrics.

\paragraph{Q: Why is the model not responsible for training/evaluation?}
\textbf{A:} Keeping training logic in the script makes it easy to compare models under the same protocol and avoids mixing I/O/metrics with differentiable code.

\paragraph{Q: What does this file export (public API)?}
\textbf{A:} The public API is the set of classes/functions other modules import (sometimes re-exported by package initializers). The key top-level definitions detected here are:

\begin{Verbatim}[fontsize=\small]
def get_args(...)  # function
def get_loaders(...)  # function
def get_gmm_mean(...)  # function
def get_gmm_label(...)  # function
def get_real_value(...)  # function
def mae_rescale_to_second(...)  # function
def test(...)  # function
\end{Verbatim}

\paragraph{Q: What are the main dependencies (imports) and why do they matter?}
\textbf{A:} Imports show whether this file is model code (torch), data code (pandas), or evaluation/analysis (sklearn). They also determine runtime requirements.

\vspace{1mm}
\VerbatimInput[firstline=1,lastline=50,fontsize=\small]{/workspace/run_d2co.py}


\paragraph{Q: What is \texttt{\detokenize{get_args}} and what responsibility does it have?}
\textbf{A:} This function is a core unit in the pipeline. Its responsibility is inferred from its name and how run scripts call it.

\paragraph{Q: What does the implementation look like for \texttt{\detokenize{get_args}}?}
\textbf{A:} Excerpt from the file:

\vspace{1mm}
\VerbatimInput[firstline=16,lastline=30,fontsize=\scriptsize]{/workspace/run_d2co.py}


\paragraph{Q: What is \texttt{\detokenize{get_loaders}} and what responsibility does it have?}
\textbf{A:} This function is a core unit in the pipeline. Its responsibility is inferred from its name and how run scripts call it.

\paragraph{Q: What does the implementation look like for \texttt{\detokenize{get_loaders}}?}
\textbf{A:} Excerpt from the file:

\vspace{1mm}
\VerbatimInput[firstline=32,lastline=38,fontsize=\scriptsize]{/workspace/run_d2co.py}


\paragraph{Q: What is \texttt{\detokenize{get_gmm_mean}} and what responsibility does it have?}
\textbf{A:} This function is a core unit in the pipeline. Its responsibility is inferred from its name and how run scripts call it.

\paragraph{Q: What does the implementation look like for \texttt{\detokenize{get_gmm_mean}}?}
\textbf{A:} Excerpt from the file:

\vspace{1mm}
\VerbatimInput[firstline=41,lastline=77,fontsize=\scriptsize]{/workspace/run_d2co.py}


\paragraph{Q: What is \texttt{\detokenize{get_gmm_label}} and what responsibility does it have?}
\textbf{A:} This function is a core unit in the pipeline. Its responsibility is inferred from its name and how run scripts call it.

\paragraph{Q: What does the implementation look like for \texttt{\detokenize{get_gmm_label}}?}
\textbf{A:} Excerpt from the file:

\vspace{1mm}
\VerbatimInput[firstline=79,lastline=83,fontsize=\scriptsize]{/workspace/run_d2co.py}


\paragraph{Q: What is \texttt{\detokenize{get_real_value}} and what responsibility does it have?}
\textbf{A:} This function is a core unit in the pipeline. Its responsibility is inferred from its name and how run scripts call it.

\paragraph{Q: What does the implementation look like for \texttt{\detokenize{get_real_value}}?}
\textbf{A:} Excerpt from the file:

\vspace{1mm}
\VerbatimInput[firstline=85,lastline=89,fontsize=\scriptsize]{/workspace/run_d2co.py}


\paragraph{Q: What is \texttt{\detokenize{mae_rescale_to_second}} and what responsibility does it have?}
\textbf{A:} This function is a core unit in the pipeline. Its responsibility is inferred from its name and how run scripts call it.

\paragraph{Q: What does the implementation look like for \texttt{\detokenize{mae_rescale_to_second}}?}
\textbf{A:} Excerpt from the file:

\vspace{1mm}
\VerbatimInput[firstline=91,lastline=99,fontsize=\scriptsize]{/workspace/run_d2co.py}


\paragraph{Q: What is \texttt{\detokenize{test}} and what responsibility does it have?}
\textbf{A:} This function is a core unit in the pipeline. Its responsibility is inferred from its name and how run scripts call it.

\paragraph{Q: What does the implementation look like for \texttt{\detokenize{test}}?}
\textbf{A:} Excerpt from the file:

\vspace{1mm}
\VerbatimInput[firstline=101,lastline=118,fontsize=\scriptsize]{/workspace/run_d2co.py}


\paragraph{Q: Example: end-to-end data flow involving this file}
\textbf{A:} Core pipeline shape (schematic):

\begin{Verbatim}[fontsize=\small]
python run_d2co.py --dataset_name kuairec --dataset_path ./dataset --device cuda:0
# Internally: load dataloaders → build model → train → evaluate
\end{Verbatim}

\paragraph{Q: Full source code of the Python file}
\textbf{A:} The complete file is included verbatim for reference.

\VerbatimInput[fontsize=\scriptsize]{/workspace/run_d2co.py}

\end{CJK*}
\end{document}